\scanpage{84}
% archive sheet number: 83
Écrivons \unsure{la relation|les relations}.
\[
\sigma_1^2=\mathrm{id}\qquad\unsure{\leftarrow}
\]
(les relations $\sigma_0^2=\sigma_2^2=\mathrm{id}$ et $\sigma_0\sigma_2=\sigma_2\sigma_0$ étant évidentes). Il suffit d'écrire
\[
\sigma_1^2(s_3)=s_3^{\unsure{*|4|' }}\qquad\text{i.e.}\qquad\sigma_1(s'_3)=s_3\qquad\text{or}
\]
\[
\begin{aligned}
\sigma_1(s'_3)&=\sigma_1s_0+\deleted{\unsure{\sigma_1e_3=s_0+\lambda}}\inserted{\lambda\sigma_1e_1+\mu\sigma_1e_2+\nu\sigma_1e_3}\\
&=s_0+\lambda e_2+\mu e_1+\nu(\lambda e_1+\mu e_2+\nu e_3)\\
&=s_0+(\mu+\nu\lambda)e_1+(\lambda+\nu\mu)e_2+\nu^2e_3
\end{aligned}
\]
\margin{$\lambda\ \mu\ \nu$}

Les conditions sont
\[
\left\{\begin{aligned}
\boxed{\begin{aligned}\nu^2&=1\\\mu&=-\nu\lambda\end{aligned}}
&\quad\begin{aligned}\deleted{\unsure{\mu+\nu\lambda=0}}\quad &(1^\circ)\\ &(2^\circ)\end{aligned}\\
\lambda=-\nu\mu\quad&\text{ou encore}\quad
\overset{\text{(compte tenu de }2^\circ\text{)}}{\lambda=\nu^2\lambda}\quad\text{qui résulte de }1^\circ,
\end{aligned}\right.
\]

Sur un corps, ces conditions se décomposent en deux possibilités
\[
\left\{\begin{aligned}\nu&=1,\ \mu=-\lambda\\\nu&=-1,\ \mu=\lambda\end{aligned}\right.
\]
% slanted
\margin{On peut \unsure{prouver} en tous les cas $\nu=1$, en écrivant $\unsure{u^3}=\mathrm{id}$ on aura $\unsure{u}s''_3=s'_3$ (où $s''_3=\sigma_0(s'_3)$) $\unsure{\ne|=}$ $\sigma_0\sigma_1s_3$ ou $\sigma_1s''_3=s'_3$, i.e. $\sigma_1\sigma_0\sigma_1(s_3)=\sigma_0\sigma_1(\unsure{s'_3|s_3})$ \unsure{---}}

Notons que
\[
\det\sigma_1=\deleted{\unsure{\det}}-\nu
\]
donc si on veut $\det\sigma_1=-1$ (\unsure{i.e.|c.à.d.|Ceci.} si on veut $\sigma_1$ une réflexion) on doit avoir $\nu=1$, $\mu=-\lambda$.

On aura alors \deleted{donc}
\[
s'_3-s_3=\deleted{\unsure{\lambda e_1+\mu e_2}}\inserted{\underbrace{(s'_3-s_0)}_{\unsure{u(e_3)}}\unsure{+|-}\underbrace{(s_3-s_0)}_{e_3}}=\lambda(e_2-e_1)=-\lambda(s_2-s_1)
\]
On \unsure{posera|pose}
\[
-\lambda=(1+\alpha)\qquad\lambda=-(1+\alpha)=\alpha'\tag*{(1)}
\]
de sorte qu'on aura
\[
s'_3-s_3=(1+\alpha)(s_2-s_1)\tag*{(2)}
\]
